% Source-aligned English for printed p.124.
\SourcePage{124}

\LayerTarget{tohoku:section:1.1:p8}
\noindent natural to choose, among all the $(A,(u_i))$ as above, a particular
system (for example by means of Hilbert's symbol $\tau$) which will be called
\textit{the product} of the family of objects $(A_i)_{i\in I}$. It is therefore
not a mere object $A$ of $\mathbf{C}$, but such an object equipped with a family $(u_i)$
of morphisms into the $A_i$, called the \textit{canonical projections of the
product} onto its factors $A_i$. The product of the $A_i$ (if it exists) is
denoted by $\displaystyle\prod_{i\in I}A_i$. If $I$ consists of a single element $i$, then
the product identifies with $A_i$ itself. We say that $\mathbf{C}$ is a
\textit{category with products} if the product of two objects of $\mathbf{C}$ always
exists (then the same is true of the product of any nonempty finite family of
objects of $\mathbf{C}$). We say that $\mathbf{C}$ is a \textit{category with infinite products}
if the product of any nonempty family of objects of $\mathbf{C}$ always exists. We have
seen that if we have two products $A=\displaystyle\prod_{i\in I}A_i$ and
$B=\displaystyle\prod_{i\in I}B_i$ corresponding to the same index set $I$, then a family
$(v_i)$ of morphisms $A_i\to B_i$ canonically defines a morphism
$v$ from $A$ to $B$, called the \textit{product of the morphisms $v_i$} and
sometimes denoted $\displaystyle\prod_{i\in I}v_i$. If the $v_i$ are monomorphisms, then so
is their product, but the analogous statement for epimorphisms is false in
general (as one sees, for example, in the category of sheaves on a fixed
topological space).

\LayerTarget{tohoku:section:1.1:p9}
Considerations dual to the preceding ones make it possible to define the notion
of a \textit{representation of an object $A$ as a sum} of a family of objects
$A_i$ by morphisms $u_i:A_i\to A$ (for every $B\in \mathbf{C}$, the natural
map
\[
  \mathrm{Hom}(A,B)\to
  \prod_{i\in I}\mathrm{Hom}(A_i,B)
\]
is bijective), of the \textit{direct sum} $\displaystyle\bigoplus_{i\in I}A_i$, equipped
with the \textit{canonical injections}
$A_i\to\displaystyle\bigoplus_{i\in I}A_i$ (which, despite their name, are not
necessarily monomorphisms), and of the \textit{sum morphism} of a family of
morphisms $u_i:A_i\to B_i$. If the $u_i$ are epimorphisms, then so
is their sum.

\LayerTarget{tohoku:section:1.2}
\noindent\textbf{1.2. Functors.}\quad
\LayerTarget{tohoku:section:1.2:p1}
Let $\mathbf{C},\mathbf{C}'$ be two categories. Recall that a \textit{covariant functor} from
$\mathbf{C}$ to $\mathbf{C}'$ is a ``function'' $F$ which associates to an object $A\in \mathbf{C}$ an
object $F(A)$ of $\mathbf{C}'$, and to a morphism $u:A\to B$ in $\mathbf{C}$ a
morphism $F(u):F(A)\to F(B)$, in such a way that
$F(I_A)=I_{F(A)}$ and $F(vu)=F(v)F(u)$. The analogous definition applies to
\textit{contravariant functors} from $\mathbf{C}$ to $\mathbf{C}'$ (which are also covariant
functors from $\mathbf{C}^{\circ}$ to $\mathbf{C}'$ or from $\mathbf{C}$ to ${\mathbf{C}'}^{\circ}$). We similarly
define functors of several variables, or \textit{multifunctors}, covariant in
some variables and contravariant in others. In general discussions, for
simplicity we shall restrict ourselves to functors of one variable. Functors
compose like functions; this composition is associative and the ``identity
functors'' play the role of identities.

\LayerTarget{tohoku:section:1.2:p2}
Let $\mathbf{C}$ and $\mathbf{C}'$ be two fixed categories, and let $F$ and $G$ be two covariant
functors from $\mathbf{C}$ to $\mathbf{C}'$. A \textit{functorial morphism} $f$ from $F$ to $G$
(also called a ``natural

\TohokuContinuation{125}
